% motivation
\chapter{Introduction and Motivation}
\label{ch:intro}

\section{Project Overview}
This project characterizes how radiation-induced bit-flip faults (Single Event Upsets, SEUs) propagate through a classical GPU-accelerated finite element solver and a GPU-simulated quantum statevector evolution of the same underlying PDE. It extends the assembly-free FEM heat solver built and optimized at MNHack7, which achieved a reduction from roughly 5\,h (naive, CPU-only) to roughly 1.78\,s (optimized, CPU+GPU) on MareNostrum 5.

\section{The Silent Data Corruption Problem}
Hardware ECC on GPUs protects DRAM. It does not protect registers, cache, shared memory, or in-transit data -- an SEU striking any of these produces a silently wrong numerical result with no crash and no error flag. This gap is the specific motivation for injecting faults below the DRAM level rather than at it: fault injection restricted to global memory alone would strike precisely the level ECC already protects, which is the wrong place to demonstrate the problem this project is actually about.

\section{Research Question}
\todo{State the one-sentence project description here once finalized. Working version from prior sync: ``SEU fault tolerance characterisation for assembly-free FEM solvers and quantum surrogates on MareNostrum 5.''}

\section{Scope and Limitations}
\todo{Pull this section together once Layers 0--3 have real results -- premature to write a scope/limitations statement before the verification layer (\autoref{ch:layer0}) has confirmed the baseline solver is even correct.}

\section{Related Work}
\label{sec:related-work}
\todo{Every entry below needs verification against a live source before this chapter is considered submission-ready. Confidence levels are an honest assessment from the last working session, not a substitute for checking -- do not cite any of these in an actual submission without confirming them first.}

\begin{itemize}
  \item Algorithm-based fault tolerance via checksums: the foundational reference. High confidence it exists roughly as cited \citep{huang1984abft}.
  \item Physics of radiation-induced soft errors in semiconductor devices: standard reference. High confidence \citep{baumann2005radiation}.
  \item Field-defining survey on silent data corruption at exascale. High confidence \citep{snir2014exascale}.
  \item Self-stabilizing / fault-tolerant iterative solvers -- the closest existing work to this project's classical control condition. Moderate confidence on exact authorship and title; finding and confirming the correct paper in this line is the single highest-priority literature task before finalizing this chapter \citep{sao2013selfstabilizing}.
  \item GPU fault-injection tooling (NVBitFI, built on NVBit; SASSIFI) -- worth citing even though hand-rolled injection is used here, for direct control over fault semantics. Moderate confidence \citep{villa2019nvbit,hari2017sassifi}.
  \item StreamFlow: explicitly excludes silent errors from its formal fault-tolerance model for hybrid workflows -- the specific, citable gap motivating the hybrid-pipeline chapter (\autoref{ch:layer8}).
  \item Standard grounding for the Kraus-operator/noise-channel treatment needed to state precisely what a quantum-side fault is (\autoref{app:fault}) \citep{nielsen2010quantum}.
\end{itemize}